\documentclass{article}
\usepackage{amsmath, amssymb, amsthm}
\usepackage{hyperref}
\usepackage{geometry}
\geometry{margin=1in}

\title{Erd\H{o}s Problem \#9}
\date{Last updated: 2025-08-31}
\author{Source: erdosproblems.com}

\begin{document}
\maketitle

\noindent\textbf{Status:} OPEN \quad \textbf{No prize}

\noindent\textbf{Tags:} number theory, additive basis, primes

\medskip
\noindent\textbf{Problem Statement:}

Let $A$ be the set of all odd integers $\geq 1$ not of the form $p+2^{k}+2^l$ (where $k,l\geq 0$ and $p$ is prime). Is the upper density of $A$ positive?

\bigskip
\noindent\textbf{Remarks:}

Crocker \cite{Cr71} proved that there are infinitely many odd integers not of this form; his proof in fact proves there are $\gg\log\log N$ such integers in $\{1,\ldots,N\}$. Pan \cite{Pa11} improved this to $\gg_\epsilon N^{1-\epsilon}$ for any $\epsilon>0$. Erd\H{o}s believed this cannot be proved by covering systems, i.e. integers of the form $p+2^k+2^l$ exist in every infinite arithmetic progression.

The sequence of such numbers is A006286 in the OEIS.

In \cite{Er80} Erd\H{o}s conjectured 'with some trepidation' that for any finite set of primes $P$ all large integers $n$ can be written as $n=m+2^k+2^l$ where $m$ is a multiple of one of the primes in $P$.

See also [10], [11], and [16].

This is discussed in problem A19 of Guy's collection \cite{Gu04}.

\bigskip
\noindent\textbf{References:}

[Cr71] Crocker, Roger, On the sum of a prime and of two powers of two. Pacific J. Math. (1971), 103-107.
 
 [Er80] Erd\H{o}s, Paul, A survey of problems in combinatorial number theory. Ann. Discrete Math. (1980), 89-115.
 
 [Gu04] Guy, Richard K., Unsolved problems in number theory. (2004), xviii+437.
 
 [Pa11] Pan, Hao, On the integers not of the form {$p+2^a+2^b$}. Acta Arith. (2011), 55-61.

\bigskip
\noindent\small{Source: \url{https://www.erdosproblems.com/9}}
\end{document}
